\chapter{The Quoting Engine}\label{ch:hf:the-quoting-engine}

A market maker changes its mind about its price tens of thousands of times a day in a busy instrument. Every change is a cancel and a new order, or
a modification; each can cross a fill already on its way, and the venue counts them all. The component that turns ``where I want to be'' into the
fewest messages that get there is the quoting engine. In this chapter's simulated market, a one-level quoter sends 758 messages an hour for 250
fills. A three-level ladder sends 980 and earns more, because its orders are already in the queue when the price moves. A hysteresis that looks
like thrift, keeping an order one tick from its target, turns the mark-out negative.

\section{From target quotes to orders}

The policies of chapters 3 to 9 produce targets: on each side, the prices and sizes the market maker wants to show. The engine owns the orders:
which are resting, which are on their way, which the venue has not yet acknowledged, which it has been asked to cancel. It compares the target with
the orders and sends the difference. Its rules (\cref{lst:hf:the-quoting-engine:diff}):

\begin{method}[Diffing a target against the working orders]\label{met:hf:the-quoting-engine:diff}
On each side, visit the working orders in price priority. An order at a target price covers it: if it holds more than the target, amend it down;
if less, add a new order for the difference, never increase it. An order at no target price is cancelled, unless it is within a hysteresis of a
target it can cover. Every target left uncovered gets a new order. Leave orders awaiting acknowledgement alone. Send cancels first, then
amendments, then new orders.
\end{method}

The order matters. Cancels go first because a cancel protects capital and a new order commits it; when messages must be rationed, the ones that
remove risk are sent. Orders in flight are left alone because the engine does not yet know whether they are in the book; acting on them twice is how
engines produce duplicate orders.

\begin{definition}[Quote ladder]\label{def:hf:the-quoting-engine:ladder}
A \emph{quote ladder}\index{quote ladder} is a set of orders resting at several consecutive prices on one side, each intended to trade: the best
level earns the spread now, and the levels behind it hold places in queues that become the best when the price moves.
\end{definition}

In a large-tick book (chapter 5) a ladder's deeper levels rarely fill where they are; their value is the place they hold. When the price ticks up,
a one-level quoter cancels its bid and joins the back of the new best queue; a ladder's second level is already in that queue.

\section{Cancel-replace, priority and in-flight states}

\begin{definition}[Cancel-replace]\label{def:hf:the-quoting-engine:replace}
A \emph{cancel-replace}\index{cancel-replace} changes a resting order's price or size in one message, the venue cancelling the old order and entering
a new one; under the usual venue rule the order keeps its place in the queue only if its size is reduced at the same price, and goes to the back for
any other change.
\end{definition}

One Quant Book 10, chapter 26 specifies the simulator's replace message with exactly this rule: a size decrease at the same price is published as a
partial cancel and keeps priority; anything else loses it. That is why the engine amends only downwards and adds a separate order to increase size.

An order has a life the engine must follow (\cref{fig:hf:the-quoting-engine:states}): sent, then acknowledged; asked to change or cancel, then acknowledged again. A fill can arrive at any
point, including after the engine has decided to cancel the order: the cancel is still travelling, or reached the venue just after the fill, which
answers ``too late to cancel''. This cancel--fill race is not an error; it is a certainty at high message rates, and the engine counts it. In the
chapter's runs a one-level quoter meets 19.5 such races an hour.

\begin{omfigure}
\begin{tikzpicture}[font=\footnotesize,>=stealth,st/.style={draw,rounded corners,fill=omThm!8,minimum width=2.3cm,minimum height=0.7cm}]
\node[st] (pn) at (0,0) {pending new};
\node[st] (lv) at (4.2,0) {live};
\node[st] (pa) at (4.2,1.8) {pending amend};
\node[st] (pc) at (8.4,0) {pending cancel};
\node[st,fill=gray!10] (gone) at (8.4,-1.8) {gone};
\draw[->] (pn) -- node[above] {ack} (lv);
\draw[->] (lv.north west) to[bend left=15] node[left] {amend down} (pa.south west);
\draw[->] (pa.south east) to[bend left=15] node[right] {ack} (lv.north east);
\draw[->] (lv) -- node[above] {cancel} (pc);
\draw[->] (pc) -- node[right] {ack} (gone);
\draw[->,dashed] (lv.south) |- node[pos=0.25,left] {fill of the whole leaves} (gone.west);
\draw[->,dashed,omDef] (pc.north east) to[out=40,in=320,looseness=5] node[right,align=left] {fill: the\\cancel--fill race} (pc.south east);
\end{tikzpicture}

\omcaption{The engine's states for one order. Solid arrows are the engine's own actions and the venue's acknowledgements; dashed arrows are fills,
which can arrive in any state; a fill received while a cancel is in flight is counted as a race. An increase in size is never an amendment: it is a
new order at the same price. Source: \texttt{firm.quoteengine}.}\label{fig:hf:the-quoting-engine:states}
\end{omfigure}

\section{Throttles, message limits and order-to-trade ratios}

\begin{definition}[Message throttle]\label{def:hf:the-quoting-engine:throttle}
A \emph{message throttle}\index{message throttle} is a limit on the rate of messages a trading system sends, imposed by the venue on each session or by
the firm on itself, and the policy for what is sent when the limit binds: which messages go first, which are deferred and which are dropped.
\end{definition}

The engine's throttle is a token bucket (One Quant Book 13, chapter 21 builds one for the gateway): tokens refill at the permitted rate up to a burst
depth, each message takes one, and when none is left the remaining actions are dropped. Dropping, not queueing, is the right policy for quotes: the
next update recomputes the difference from the latest target, so a stale action is never sent late.

Venues limit messages in three ways: a hard rate per session (a throttle rejects the excess), a charge on messages beyond a ratio to trades, and the
regulatory \emph{order-to-trade ratio} of One Quant Book 10, chapter 24.

\begin{dated}{2026-09}{Message limits in rules and fee schedules}\label{dat:hf:the-quoting-engine:limits}
Under Commission Delegated Regulation (EU) 2017/566, each trading venue computes, for each member and instrument at least at the end of every
session, the ratio of unexecuted orders to transactions in volume terms and in number terms (orders over transactions, minus one), counting each order
type as its annex prescribes, and a member exceeding the venue's maximum ratio in either measure has exceeded it. CME Group's messaging efficiency
programme, as filed with the CFTC in 2019, weights messages by type into a messaging score, divides it by the firm's traded volume in a product group
during regular trading hours, and compares the resulting volume ratio with product-group benchmarks set each quarter.
\end{dated}

\section{Self-trade prevention and multiple strategies}

A firm running several strategies in one instrument can find its own bid crossing its own offer: a quoting strategy resting at the ask and a
momentum strategy buying. The trade is real on the venue's books but moves no risk, pays two fees, and prints a wash trade that surveillance will
see (One Quant Book 9, chapter 29). Venues offer \emph{self-trade prevention} (One Quant Book 10, chapter 26): orders tagged with the same group are
never matched against each other, and the venue cancels the newer, the older or both, as the firm chose. The engine's side of the problem is to
know every strategy's resting orders and to route internal crossings to the netting logic of chapter 11 before they reach the venue.

\section{Quote ladders and the line not to cross}

A ladder of orders on one side, cancelled and moved as the market moves, looks from the outside like the layering of One Quant Book 9, chapter
29: orders at several prices that rarely trade. The difference is intent. The statute under which the Seventh Circuit upheld Coscia's conviction in
2017 defines spoofing as ``bidding or offering with the intent to cancel the bid or offer before execution''. The court rejected his argument that
the definition was vague because high-frequency traders cancel 98\% of their orders. Cancelling most orders is normal; placing them never meaning to
let them trade is not. A market maker's ladder is defensible when each level is priced and sized to trade, fills are taken when they come, both sides
are quoted, and the cancellation logic is the policy's, documented. The engine is where that evidence is produced: it logs, for every order, the
target it served.

\section{Tutorial: one engine, four configurations}\label{tut:hf:the-quoting-engine:tut}

\begin{tutorial}
\textbf{Goal.} Drive the engine from a quoting policy on the simulated market and measure messages, fills, races, queue places and edge.
\textbf{End state:} the table below.
\begin{tutsteps}
\item \textbf{The diff} (\cref{met:hf:the-quoting-engine:diff}), in \texttt{firm.quoteengine.QuoteEngine.update}; its C++20 and Rust twins produce the
same 1,154 actions on a 3,000-event fixture.
\omcode{code/firm/quoteengine/firm_quoteengine.py}{71}{102}{Cover targets with working orders, amend down, add, cancel, or keep within the hysteresis.}\label{lst:hf:the-quoting-engine:diff}
\item \textbf{Issuing under a budget}: cancels first, then amendments, then new orders, one token each.
\omcode{code/firm/quoteengine/firm_quoteengine.py}{103}{124}{The token bucket decides what is sent; the rest is re-derived at the next update.}\label{lst:hf:the-quoting-engine:issue}
\item \textbf{The quoter.} \texttt{hf\_engine.EngineQuoter} targets the others' best prices (and the prices behind them for a ladder), one lot per
level, and translates the engine's actions into orders; on \texttt{firm.tape}, which has no modify message, an amendment becomes a cancel and a new
order.
\item \textbf{Compare} four configurations on four twenty-minute sessions with \texttt{compare()}.
\end{tutsteps}
\textbf{What to change next.} Run the same quoter on \texttt{firm.exchsim}, whose replace message keeps priority on a size decrease; add a second
strategy in the same instrument and count the self-trades prevented.
\end{tutorial}

\begin{center}
\footnotesize
\begin{tabular}{@{}lrrrrrr@{}}
\toprule
configuration & messages & fills & races & 10-s mark-out & edge & median age\\
& an hour & an hour & an hour & (ticks a share) & (\$ an hour) & of a fill (s)\\
\midrule
one level & 758 & 250 & 19.5 & 0.141 & 35.25 & 9.6\\
three levels & 980 & 282 & 20.3 & 0.144 & 40.50 & 17.3\\
three levels, hysteresis of two ticks & 806 & 364 & 1.5 & $-0.064$ & $-23.25$ & 9.6\\
three levels, 2 messages a second & 884 & 282 & 9.8 & 0.154 & 43.50 & 17.7\\
\bottomrule
\end{tabular}
\end{center}

Four lessons (\cref{fig:hf:the-quoting-engine:configs}). The ladder sends 29\% more messages and earns 15\% more edge an hour, because its fills come from orders that waited in the queue (a
median of 17 seconds against 10). The hysteresis saves messages but keeps orders a tick away from where the policy wants them, including a bid one
tick above a best bid that has fallen: that order is alone at a stale price and gets picked off, and the mark-out turns negative; hysteresis in a
one-tick book must only ever keep orders behind the target, never in front. The throttle of two messages a second drops 1,396 derived actions an hour
and loses nothing: the actions dropped were re-derived and sent when a token was free, and the budget cut the races by half. The messages per fill,
3.0 to 3.5, are what a venue's order-to-trade ratio would count.

\begin{omfigure}
\begin{tikzpicture}
\begin{axis}[width=11cm,height=5.4cm,ybar,bar width=10pt,ymin=0,ymax=1100,xmin=-0.5,xmax=3.5,xtick={0,1,2,3},
  xticklabels={one level,three levels,hysteresis,throttle},ylabel={\small an hour},tick label style={font=\footnotesize},
  axis y line*=left,area legend,table/col sep=comma,legend style={font=\scriptsize,at={(0.5,-0.22)},anchor=north},legend columns=3]
\addplot[fill=omThm!60,draw=omThm] table[x=i,y=messages]{figdata/hft/10-the-quoting-engine/configs.csv};
\addplot[fill=gray!40,draw=gray] table[x=i,y=fills]{figdata/hft/10-the-quoting-engine/configs.csv};
\addlegendentry{messages}\addlegendentry{fills}
\addlegendimage{line legend,omDef,only marks,mark=*}\addlegendentry{edge, \$ an hour (right)}
\end{axis}
\begin{axis}[width=11cm,height=5.4cm,ymin=-40,ymax=60,xmin=-0.5,xmax=3.5,axis y line*=right,axis x line=none,
  ylabel={\small edge (\$ an hour)},tick label style={font=\footnotesize},table/col sep=comma]
\addplot[omDef,only marks,mark=*,mark size=2pt] table[x=i,y=edge]{figdata/hft/10-the-quoting-engine/configs.csv};
\draw[omDef!50,dotted] (axis cs:-0.5,0) -- (axis cs:3.5,0);
\end{axis}
\end{tikzpicture}

\omcaption{Messages and fills an hour (bars, left) and the ten-second edge in dollars an hour (dots, right) for the four configurations of the
tutorial: one level, a three-level ladder, the ladder with a two-tick hysteresis, and the ladder under a throttle of two messages a second with a
burst of five; one lot a level, four twenty-minute sessions. Data: \texttt{hf\_engine.compare}.}\label{fig:hf:the-quoting-engine:configs}
\end{omfigure}

\section{Build: the quoting engine}\label{bld:hf:the-quoting-engine:quoteengine}

\begin{build}
\bfield{Purpose} Turn target quotes into the fewest messages under a venue's priority rules and a message budget, and track every order's state.
\bfield{Interface} \texttt{QuoteEngine(min\_move, min\_size, rate, burst)}: \texttt{update(t, side, targets) -> actions} (\texttt{('new', oid, side, price,
qty)}, \texttt{('amend', oid, qty)}, \texttt{('cancel', oid)}); \texttt{ack(oid)}, \texttt{fill(oid, qty)}, \texttt{orders}, \texttt{stats}
(new, amend, cancel, dropped, fills, races); \texttt{TokenBucket(rate, burst).take(t)}. C++20: \texttt{cpp/firm\_quoteengine.hpp}; Rust:
\texttt{rust/}.
\bfield{Rules} Amend only downwards; orders in flight are untouched; cancels before amendments before new orders; an empty bucket drops actions,
never queues them; a fill during a pending cancel is counted as a race.
\bfield{Acceptance tests} \texttt{code/firm/quoteengine/tests/}: no message when the book already matches the target; a price move cancels one
order and adds one; a size decrease amends, an increase adds; the hysteresis keeps an order within one tick; the throttle drops and the next update
re-derives; a cancel--fill race is counted; the Python, C++20 and Rust engines produce the same action log on the fixture.
\bfield{Stretch} Replace messages with the venue's priority rule on \texttt{firm.exchsim}; per-strategy self-trade groups; an order-to-trade budget per
instrument and session.
\end{build}

\begin{omsources}
\item Commission Delegated Regulation (EU) 2017/566 on the ratio of unexecuted orders to transactions.
\item CME Group, CME Globex Messaging Efficiency Program, rule filing with the CFTC, November 2019.
\item \emph{United States v. Coscia}, 866 F.3d 782 (7th Cir. 2017).
\end{omsources}

\section{Exercises}

\begin{exercise}[$\star$]\label{exo:hf:the-quoting-engine:1}
The engine rests 300 shares at 100 and the target becomes 100 shares at 100. What does it send, and does the order keep its place?
\end{exercise}

\begin{exercise}[$\star$]\label{exo:hf:the-quoting-engine:2}
A member sent 12,000 orders and made 400 transactions in an instrument in a session. What is its ratio of unexecuted orders to transactions in number
terms?
\end{exercise}

\begin{exercise}[$\star$]\label{exo:hf:the-quoting-engine:3}
A token bucket refills at 2 tokens a second with a depth of 5. It is full at $t=0$ and five actions are sent at once. How many can be sent at
$t=1.2$?
\end{exercise}

\begin{exercise}[$\star\star$]\label{exo:hf:the-quoting-engine:4}
Why should a throttled engine drop actions rather than queue them?
\end{exercise}

\begin{exercise}[$\star\star$]\label{exo:hf:the-quoting-engine:5}
Why did the two-tick hysteresis turn the mark-out negative, and how would you fix it without losing its message savings?
\end{exercise}

\begin{exercise}[$\star\star$]\label{exo:hf:the-quoting-engine:6}
What makes a \omterm{def:hf:the-quoting-engine:ladder}{quote ladder} defensible against an accusation of layering?
\end{exercise}

\begin{exercise}[$\star\star\star$]\label{exo:hf:the-quoting-engine:7}
\emph{Coding.} Replay the component's fixture through an engine with \texttt{rate=5} and \texttt{burst=2}
(\texttt{hf\_engine.fixture\_stats}), report the counts, and explain why they say nothing about what a tighter throttle would do.
\end{exercise}

\begin{exercise}[$\star\star\star$]\label{exo:hf:the-quoting-engine:8}
\emph{Find the flaw.} ``To save messages we increase an order's size with a modify instead of sending a second order.''
\end{exercise}

\section{Problem: Ten Thousand Changes of Mind}

\begin{problem}[{Weekend problem --- ten thousand changes of mind}]\label{pb:hf:the-quoting-engine:1}
A market maker's engine must keep a three-level ladder current in a busy one-tick stock without exceeding its message budget or its venue's
order-to-trade ratio.

\medskip\noindent\textbf{Part I --- The engine.}
\begin{enumerate}
\item State the diff of \cref{met:hf:the-quoting-engine:diff} and justify the order in which actions are sent.
\item Define a \omterm{def:hf:the-quoting-engine:ladder}{quote ladder} and say what its deeper levels are for in a one-tick book.
\item Define a \omterm{def:hf:the-quoting-engine:replace}{cancel-replace} and the priority rule, and explain why the engine never amends upwards.
\item What is a cancel--fill race, and how often did it happen here?
\end{enumerate}

\medskip\noindent\textbf{Part II --- Limits.}
\begin{enumerate}[resume]
\item Define a \omterm{def:hf:the-quoting-engine:throttle}{message throttle} and describe the token bucket.
\item Summarise the two message limits of the dated box.
\item What is self-trade prevention and why does a multi-strategy firm need it?
\item Where is the line between a ladder and layering, and what did the Coscia court say about cancellation rates?
\end{enumerate}

\medskip\noindent\textbf{Part III --- Measurements.}
\begin{enumerate}[resume]
\item Give messages, fills and edge for the one-level and three-level quoters.
\item Why do the ladder's fills come from older orders?
\item What did the two-tick hysteresis do, and why?
\item What did the throttle drop, and what did it cost?
\end{enumerate}

\medskip\noindent\textbf{Part IV --- The verdict.}
\begin{enumerate}[resume]
\item State the \emph{named result}: messages an hour and the median age of a fill under the one-level, three-level and throttled configurations,
against the edge each earns.
\item Which configuration would you run, and with what change to the hysteresis?
\item What would the messages-per-fill figures mean for a venue's order-to-trade ratio?
\item How would results change on a venue whose replace keeps priority on size decreases?
\item Why must the Python, C++ and Rust engines agree action for action?
\item What does the engine log to defend the firm's quoting?
\item How would a second strategy in the same stock change the engine?
\item In one sentence: what does a quoting engine optimise?
\end{enumerate}
\end{problem}

\section{Interview questions}

\begin{interviewq}[$\star$ \iqroles{developer}]\label{iq:hf:the-quoting-engine:1}
Your engine sends a cancel and immediately receives a fill for the same order. What states must it handle, and what does it tell the strategy?
\end{interviewq}

\begin{interviewq}[$\star\star$ \iqroles{developer}]\label{iq:hf:the-quoting-engine:2}
Design the data structure for working orders so that the diff against a ten-level ladder runs in constant time per level.
\end{interviewq}

\begin{interviewq}[$\star\star$ \iqroles{trader}]\label{iq:hf:the-quoting-engine:3}
Your venue caps you at 50 messages a second and the market is moving fast. Which messages do you send first?
\end{interviewq}

\begin{interviewq}[$\star\star$ \iqroles{risk}]\label{iq:hf:the-quoting-engine:4}
A regulator asks why your firm cancels 97\% of its orders. What do you show them?
\end{interviewq}

\begin{interviewq}[$\star\star$ \iqroles{researcher}]\label{iq:hf:the-quoting-engine:5}
How would you measure the value of a place in the queue that your ladder preserves when the price moves?
\end{interviewq}

\begin{interviewq}[$\star\star\star$ \iqroles{developer}]\label{iq:hf:the-quoting-engine:6}
Prove that a token bucket of rate $r$ and depth $b$ never lets more than $b+r\,T$ messages through in any window of length $T$.
\end{interviewq}
